% Define document class
\documentclass[twocolumn]{aastex631}
\usepackage{showyourwork}
\usepackage{amsmath}
\usepackage{bm}
\usepackage{listings}
\usepackage{natbib}
\usepackage{xcolor}


\definecolor{codegreen}{rgb}{0,0.6,0}
\definecolor{codegray}{rgb}{0.5,0.5,0.5}
\definecolor{codepurple}{rgb}{0.58,0,0.82}
\definecolor{backcolour}{rgb}{0.95,0.95,0.92}
\lstset{
    backgroundcolor=\color{white},   
    commentstyle=\color{codegreen},
    keywordstyle=\color{magenta},
    numberstyle=\tiny\color{codegray},
    stringstyle=\color{codepurple},
    basicstyle=\ttfamily\footnotesize,
    breakatwhitespace=false,
    breaklines=true,
    captionpos=b,
    keepspaces=true,
    numbers=left,
    numbersep=5pt,
    showspaces=false,
    showstringspaces=false,
    showtabs=false,
    tabsize=1
}

\definecolor{tedcommentcolor}{HTML}{e17701}
\newcommand{\TJ}[1]{\textcolor{tedcommentcolor}{#1}}

\DeclareMathOperator{\Span}{span}

\defcitealias{stevenson2020}{S20}


% Begin!
\begin{document}

% Title
\title{The Variable Planetary Infrared Excess Method}

% Author list
\author{Ted Johnson}
\affiliation{UNLV}
\author{Avi Mandell}
\affiliation{GSFC}
\author{Rebecca Martin}
\affiliation{UNLV}

% Abstract with filler text
\begin{abstract}
    The Planetary Infrared Excess method to measure the thermal emission of an exoplanet is analogous to
    detecting the infrared excess of a circumstellar disk or low-mass companion -- except that the 
    exceptionally low contrast between planet and star require extreme precision and make the method
    easily susceptible to bias. We propose a new method to analyze PIE time-series data that naturally
    accounts for stellar variability and is agnostic to stellar spectral models. This method treats
    the observation as a vector space and isolates variability present at long wavelengths that is
    uncorrelated to short-wavelength variability.
\end{abstract}

% Main body with filler text
\section{Introduction}
\label{sec:intro}
Two decades ago, the Spitzer Space Telescope detected photons emitted by an exoplanet for the first time, paving the way for the modern study of exoplanets in thermal emission \citep{deming2005}. In the years that followed, secondary eclipse measurements -- determination of a planet's flux found by the dimming of a system as the planet passes behind its host star -- became a standard way to characterize exoplanets with Spitzer \citep[e.g.][]{kreidberg2019} and more recently with the James Webb Space Telescope \citep[JWST, e.g.][]{greene2023,kempton2023,weinermansfield2024}.

\begin{figure*}
    \centering
    \includegraphics[width=0.8\textwidth]{figures/nearby_mdwarfs.pdf}
    \caption{Known potentially-rocky exoplanets within 20 pc. The vast majority of the closest planets are non-transiting and their atmospheres cannot be characterized with currently-available techniques. Data from the NASA Exoplanet Archive, accessed on \variable{output/nearby_mdwarfs_date.txt}. An interactive version of this figure is available through the NASA Exoplanet Modeling and Analysis Center \TJ{Talk to Joe about getting this working.}}
    \script{nearby_mdwarfs.py}
    \label{fig:nearby_mdwarfs}
\end{figure*}

The eclipse method, however, relies on the observer's line-of-sight being in the orbital plane of the planet, a probability that is approximately $R_*/a$ (0.5\% for Earth, $\sim 1\%$ for potentially-habitable planets in M-dwarf systems). Figure \ref{fig:nearby_mdwarfs} shows us that indeed the population of known nearby planets is mostly non-transiting. Within 10pc there are no cool transiting planets, leaving Gliese-12b and TRAPPIST-1e as the best candidates to characterize atmospheres in the habitable zone (HZ), despite predictions that sufficient signal-to-noise would require hundreds of eclipses and be infeasible in the lifetime of JWST \citep{morley2017}. This leaves no proven method to characterize {\em any} HZ rocky planet using thermal emission.

The future flagship Habitable Worlds Observatory (HWO) will have the ability to spatially resolve Earth-like planets around Sun-like stars and characterize atmospheres with reflected light, but will be severely limited by the exposure time required for each target. Optimizing all components of the planned survey has become a field in itself \citep[e.g.][]{latouf2023,tuchow2024}, but as Earth is the only known habitable planet, instrument optimization studies are limited to detections of Earth's atmosphere at various epochs throughout geologic history \citep[e.g. Modern vs. Archean Earth, see][]{latouf2024}. Characterization of additional HZ planetary atmospheres (even if deemed uninhabitable) would give critical insight into the population that HWO will explore.

ESA's Large Interferometer For Exoplanets (LIFE) mission will likely fly concurrently with HWO. The science output from the two observatories will complement each other well, but neither will inform the other's instrument design. Hence, it is necessary for the exoplanet community to prioritize characterizing some HZ planets in the decades prior to HWO and LIFE.

\citet{stevenson2020} introduced a new method to measure the thermal emission of an exoplanet that does not require an eclipse. Called Planetary Infrared Excess (PIE), this technique is akin to the way that circumstellar disks can be detected by looking for excess infrared light emitted by a stellar system. The excess, in the PIE case, is very small -- on the order of $10^{-5}$ with respect to the star, and so much care is needed to successfully characterize at atmosphere this way. A major challenge of the PIE method is that it requires a phyiscs-based stellar model to be accurate across a broad wavelength range. Recent JWST studies have found that even the best models can be reliable only to within $\sim 1\%$ \citep{moran2023}. Additionally, it has been found that inaccuracies in stellar models lead to significant systematic errors in the absolute flux estimates of eclipsing planets and the worst inaccuracies are in models of M-dwarfs, which host most nearby planets \citep{fauchez2025}.

In this study, we introduce the Variable Planetary Infrared Excess (VPIE) method, a technique that extends the concepts of the PIE method to separate planetary and stellar flux variability without the need to a physics-based stellar model. We describe VPIE in Section \ref{sec:vpie}.

\section{The VPIE Method}
\label{sec:vpie}

VPIE is a PIE time-series analysis method that aims to measure the variability of a planet as it orbits is host. 


The VPIE (Variable Planetary Infrared Excess) method aims to overcome the challenges described in Section \ref{sec:pie} by making a concession: we will only measure changes in the planet's flux. Just like in Section \ref{sec:pie}, we will start with an observation $\mathbf{F}$ with associated uncertainty $\mathbf{\Sigma}$, except that each of these is an $m \times n$ matrix where $m$ is the number of epochs and $n$ is the number of wavelength channels. We also define wavelength $\bm{\lambda}$ and time $\bm{t}$ such that ${\rm F}_{ij}$ is the flux at time $t_i$ and wavelength $\lambda_j$.

Our goal is the same: to compute some stellar model $\tilde{\mathbf{F}}$ to subtract from $\mathbf{F}$ in order to recover information from the planet. Again we do our fitting at short wavelengths only:
\begin{equation}
    \mathbf{F}_\text{NIR} - \tilde{\mathbf{F}}_\text{NIR} = \mathbf{E}_\text{NIR}\, ,
\end{equation}
where $\mathbf{E}$ is an error term defined as
\begin{equation}
    \label{eq:def_E}
    \mathbf{E} \equiv \mathbf{F}_\text{s} - \tilde{\mathbf{F}}_\text{s}\, .
\end{equation}

Rather than compute $\tilde{\mathbf{F}}$ from some physical model, we will approximate as a matrix whose rows are linear combinations of a small subset of the rows of $\mathbf{F}$. We can choose a set of $q$ integers $s$, whose values are the indices of this ``spectral basis''. Let $\mathbf{B}$ be a $q\times n$ matrix
\begin{equation}
    \label{eq:def_B}
    B_{kj} = F_{s_kj} \, \text{for}\, k=1,\,2,\,...,\,q\, ,
\end{equation}
i.e. the $k$th row of $\mathbf{B}$ is the $s_k$th row of $\mathbf{F}$.

We can also define a coefficient matrix $\mathbf{A}$, which has shape $m\times q$. The model observation can be written
\begin{equation}
    \tilde{\mathbf{F}} = \mathbf{A}\mathbf{B}\, .
\end{equation}

Choosing the best values for $\mathbf{A}$ and $s$ is discussed in Sections \ref{subsec_detA} and \ref{subsec:choose_B}, respectively. However, suppose we had found some optimal set of bases, described by $\hat{s}$, and then calculated $\mathbf{A}$ so that $\mathbf{E}_\text{NIR}$ is minimized. When computing the residual of the full wavelength range, we will introduce the matrix $\mathbf{\Delta}$, which is an additional error term which we attribute to the planet.
\begin{equation}
    \label{eq:measurable}
    \mathbf{F} - \tilde{\mathbf{F}} = \mathbf{E} + \mathbf{\Delta}\, .
\end{equation}
As $\mathbf{F}$ is measurable and $\tilde{\mathbf{F}}$ can be calculated from it, the quantity $\mathbf{E} + \mathbf{\Delta}$ is measurable. We can decompose $\mathbf{F}$ and $\tilde{\mathbf{F}}$ to find its meaning. Recall
\begin{equation}
    \mathbf{F} = \mathbf{F}_\text{s} + \mathbf{F}_\text{p}\, ,
\end{equation}
so
\begin{equation}
    \mathbf{B} = \mathbf{B}_\text{s} + \mathbf{B}_\text{p}\, ,
\end{equation}
and $\tilde{\mathbf{F}}$ can be written
\begin{equation}
    \begin{array}{rccccc}
        \tilde{\mathbf{F}} & = & \mathbf{A}&(\mathbf{B}_\text{s} &+& \mathbf{B}_\text{p})\, , \\
        & = & & \tilde{\mathbf{F}}_\text{s} &+& \tilde{\mathbf{F}}_\text{p}\, .
    \end{array}
\end{equation}
We can then rewrite Equation (\ref{eq:measurable}):
\begin{equation}
    \mathbf{E} + \mathbf{\Delta} = \mathbf{F}_\text{s} - \tilde{\mathbf{F}}_\text{s} + \mathbf{F}_\text{p} - \tilde{\mathbf{F}}_\text{p}\, .
\end{equation}

Given the definition of $\mathbf{E}$ in Equation (\ref{eq:def_E}), we see that
\begin{equation}
    \mathbf{\Delta} =  \mathbf{F}_\text{p} - \tilde{\mathbf{F}}_\text{p}\, .
\end{equation}

If $\mathbf{E}$ is small, then $\mathbf{F} - \tilde{\mathbf{F}}\approx \mathbf{\Delta}$, and we have measured some signal from the planet. Importantly, the expected value of $\mathbf{\Delta}$ can be calculated for any phase curve model, and therefore we can use our measured $\mathbf{\Delta}$ to infer properties of the planet.

In contrast to the method described in Section \ref{sec:pie}, our stellar error term $\mathbf{E}$ is not dominated by the limitations of a physics-based stellar model. Instead, there are two main sources of error: 1) complex variability: the time-dependent behavior of the star is too complicated to be represented by the $q$ basis vectors described by $\hat{s}$, and 2) measurement error, which in practice is mostly photon noise.

\TJ{(\#1 is unknown without real dataset, \#2 very doable)}


\subsection{Determination of $\mathbf{A}$}
\label{subsec_detA}
Given some basis matrix $\mathbf{B}$ that is based on a set of indices $s$, we need to determine the coefficient matrix $\mathbf{A}$. The best choice is one that minimizes the normalized residual matrix $\mathbf{E}$ at short wavelengths.

The coefficients $A_{ik}$ that correspond to a single epoch $t_i$ are straight forward to compute using least-squares. We are looking at a single row of each of our data matrices, and so we will write them as vectors:
\begin{equation}
    \bm{\epsilon}^\top = \bm{f}^\top - \mathbf{B}\bm{a}^\top\, .
\end{equation}

We can also define a diagonal weight matrix using the uncertainties:
\begin{equation}
    {W}_{ij} = \frac{\delta_{ij}}{\sigma_i \sigma_j}\, .
\end{equation}

We would like to use least-squares to solve the equation
\begin{equation}
    \mathbf{W}\bm{f}^\top = \mathbf{W}\mathbf{B}\bm{a}^\top\, .
\end{equation}
for $\bm{a}$, which has the solution
\begin{equation}
    \label{eq:a_transpose}
    \bm{a}^\top = (
        \mathbf{B}\mathbf{W}^2\mathbf{B}^\top
    )^{-1} \mathbf{B}\mathbf{W}^2 \bm{f}^\top\, .
\end{equation}

Here we have computed $\bm{a}$, which is a row of the $m \times q$ matrix $\mathbf{A}$. It is important to point out that, while $\mathbf{B}$ is uniquely determined by the choice of $s$ for a given problem, we have only computed $\mathbf{W}$ for a given value of the time index $i$. To compute $\mathbf{A}$ we must evaluate Equation (\ref{eq:a_transpose}) $m$ times:
\begin{equation}
    \label{eq:calc_A}
    \mathbf{A} = \begin{bmatrix}
        \bm{a}_1 \\
        \bm{a}_2 \\
        \vdots \\
        \bm{a}_m
    \end{bmatrix}\ ,
\end{equation}
where the subscript $i$ denotes that $\bm{a}$ is evaluated considering the $i$th time index. That is
\begin{equation}
    \label{eq:a_func}
    \bm{a}_i^\top = (
        \mathbf{B}\mathbf{W}_i^2\mathbf{B}^\top
    )^{-1} \mathbf{B}\mathbf{W}_i^2\bm{f}_i^\top\, ,
\end{equation}
where $\mathbf{B}$ is defined in Equation (\ref{eq:def_B}), $\bm{f}_i$ is the $i$th row of $\mathbf{F}$, and $\mathbf{W}_i$ is defined
\begin{equation}
    (\mathbf{W}_i)_{jk} = \frac{\delta_{jk}}{\Sigma_{ij}\Sigma_{ik}}\, .
\end{equation}

The above relations allow $\mathbf{A}$ to be determined given $\mathbf{F}$, $\mathbf{\Sigma}$, and $s$.

\subsubsection{Time-averaged uncertainties}
In some cases it may be appropriate to replace $\mathbf{\Sigma}$ with the time-averaged uncertainty $\langle \mathbf{\Sigma} \rangle$. This is justified if we consider that, for photon noise, $\Sigma_{ij} \propto \sqrt{F_{ij}}$. A variable star may mean a single wavelength channel changes by several percent over the course of an observation, but $\mathbf{F}$ may vary by several orders of magnitude as a function of wavelength. Averaging $\mathbf{\Sigma}$ significantly reduces computational cost, as otherwise Equation (\ref{eq:calc_A}) requires inverting a $q\times q$ matrix $m$ times. Now, Equation (\ref{eq:a_func}) is written
\begin{equation}
    \label{eq:a_func_mean}
    \bm{a}_i^\top = (
        \mathbf{B}\langle\mathbf{W}\rangle^2\mathbf{B}^\top
    )^{-1} \mathbf{B}\langle\mathbf{W}\rangle^2\bm{f}_i^\top\, ,
\end{equation}
and the matrix $ (\mathbf{B}\langle\mathbf{W}\rangle^2\mathbf{B}^\top)^{-1} \mathbf{B}\langle\mathbf{W}\rangle^2$ can be cached after the first computation.


\subsection{Choosing the basis $\mathbf{B}$}
\label{subsec:choose_B}

Our goal is to find the optimal $\mathbf{B}$, which is defined by the optimized set of indices $\hat{s}$. To do this we will test sets $s$ and compare them using the Bayes Information Criterion (BIC, \TJ{ref?}), a model selection criterion for which smaller values are preferred:
\begin{equation}
    \text{BIC} = q\ln{(mn)} - 2\ln{(\hat{L})}\, ,
\end{equation}
where $\hat{L}$ is the maximized likelihood function. This is the probability of observing $\mathbf{F}$ given the prediction given by reconstruction $\tilde{\mathbf{F}}$ and measurement uncertainties $\mathbf{\Sigma}$.

The maximized likelihood is
\begin{equation}
    \hat{L} = \prod_{i=1}^m \prod_{j=1}^n \frac{1}{\sqrt{2\pi \Sigma_{ij}^2}} \exp{\left (-\frac{(F_{ij} - \tilde{F}_{ij})^2}{2\Sigma_{ij}^2}\right )}\, .
\end{equation}


$\tilde{\mathbf{F}} = \mathbf{A}\mathbf{B}$ has already been computed using Equation (\ref{eq:calc_A}), so this likelihood is maximized. To compute the BIC, we need $\ln{\hat{L}}$:
\begin{equation}
        \ln{\hat{L}} = \sum_{i=1}^m \sum_{j=1}^n -\frac{1}{2} \left (\ln{(2\pi\Sigma_{ij}^2)} + \frac{(F_{ij} - \tilde{F}_{ij})^2}{\Sigma_{ij}^2}\right )\, .
\end{equation}

The BIC is then written
\begin{equation}
    \label{eq:bic2}
    \text{BIC} = q\ln{(mn)} + \sum_{ij} \left (\ln{(2\pi\Sigma_{ij}^2)} + \frac{(F_{ij} - \tilde{F}_{ij})^2}{\Sigma_{ij}^2}\right )\, .
\end{equation}

The only values in Equation (\ref{eq:bic2}) that change for any given $s$ are $q$, the number of elements in $s$, and $\mathbf{\tilde{F}}$, the reconstructed spectrum. Therefore, there exists some function $\text{BIC}(s)$ that can be used to compare any choice $s$ to another choice $s'$.

All that is left now is to test various values of $s$, and choose the one for which BIC is minimized. There are many ways to do this. The simplest is to start at low values of $q$ and step through each possible value of $s$ sequentially, finding a local minimum for each value of $q$. When increasing $q$ no longer leads to decreased BIC, the search is stopped and the previous minimum is returned as the solution.

\begin{lstlisting}[
    label={ls:soln1}, caption={Minimizing BIC though an exhaustive search at low $q$}
]
minima = {}
for q in 0..m-1:
    sets = combos(0..m-1, q)
    vals = []
    for s in sets:
        vals.append(bic(s))
    imin = argmin(vals)
    if q > 1:
        if vals[imin] < minima[q-1]["val"]:
            minima[q] = {
                "val":vals[imin],
                "s":sets[imin]
            }
        else:
            return minima[q-1]["s"]
    else:
            minima[q] = {
                "val":vals[imin],
                "s":sets[imin]
            }
\end{lstlisting}

This method, shown as pseudocode in Listing \ref{ls:soln1} is exhaustive, but quickly becomes computationally expensive. There are $m\choose q$ combinations for each loop to consider, so the time complexity of each loop iteration is $\mathcal{O}(m^q)$.

A less expensive method is to find local minimum with $q=1$, and then add a single index at a time as $q$ is increased. This will not necessarily find the global minimum of the BIC, but searches higher values of $q$ much faster as the time complexity is $\mathcal{O}(m)$. A pseudocode implementation is shown in Listing \ref{ls:sequential}.

\begin{lstlisting}[
    label={ls:sequential},caption={Minimizing BIC by searching for the next best index for each value of $q$}
]
s = []
vals = []
for _ in 1..m-1:
    _is = []
    _vals = []
    for i in 1..m-1 if i not in s:
        _s = s + [i]
        _is.append(i)
        _vals.append(bic(_s))
    imin = argmin(_vals)
    s.append(_is[imin])
    vals.append(_vals[imin])
q_best = argmin(vals) + 1
s = s[:q_best-1]
return s
\end{lstlisting}



\bibliography{syw,pie}

\end{document}
